\section{Introduction}
\label{sec:introduction}

Federated learning (FL) trains a shared model across decentralised
clients without centralising their data~\cite{mcmahan2017fedavg}. In
the consumer-IoT setting (home assistants, smart-home sensors,
wearable biometrics, connected vehicles) this is existential:
participants are unwilling, and increasingly not legally permitted,
to upload raw data. The FL literature has accordingly focused on
\emph{technical} privacy (secure aggregation, leakage bounds,
poisoning robustness) while a second, equally binding axis remains
under-addressed: \emph{regulatory} privacy. GDPR Article~7 requires
consent to be demonstrably given, purpose-granular, and freely
withdrawable; Article~17 grants erasure, which, applied to a
trained model, means the model must \emph{forget} a revoked
participant's contribution. The EU AI Act~\cite{aiact2024}
(Arts.~15/26) and CCPA~\cite{ccpa2018} \S~1798.105 impose analogous,
auditable obligations. None of these regimes assumes the deployment
\emph{chooses} to comply; each demands demonstrable, audit-grade
evidence.

Three research threads each address one corner of this problem but
not the whole. Consent-management ledgers (MedRec
\cite{azaria2016medrec}) record consent immutably but cannot gate
downstream training. Machine-unlearning methods (SISA
\cite{bourtoule2021sisa}, FedEraser~\cite{liu2021federaser}) remove a
participant's contribution but leave the trigger and verification
off-chain. Blockchain-FL systems~\cite{kim2020blockchained,li2020bflc}
log model updates on-chain without consent semantics or
forgetting verification. No published work binds the three.

We frame the gap as a missing \emph{edge}, not a missing vertex.
Consent must stay authoritative over which data enters which round;
unlearning (by any method) must be \emph{verifiable} rather than
self-asserted; and verification, against membership inference (the
canonical operational test for
forgetting~\cite{shokri2017mia,carlini2022lira,choquettechoo2021labelonly}),
must in turn drive the consent state, closing the loop. Prior
work implements the vertices; \sysname supplies the edges. Enforcing
them off-chain would merely move trust to the data controller;
putting the loop on a public smart-contract substrate removes that
trust but introduces a cost question: a single revocation on
Ethereum mainnet costs dollars, prohibitive at IoT scale. We show
the identical protocol on a public, gas-free, EVM-compatible chain
with PoA$^2$ consensus costs \emph{exactly} \$0 per operation.
This work make the following key contributions:
\begin{enumerate}[leftmargin=*]
    \item A \textbf{Consent Lifecycle State Machine} (\clsm): six
        states, nine transitions, four formally-stated safety and
        liveness properties, instantiated isomorphically as a
        soulbound ERC-721 contract.
    \item A \textbf{Verifiable Unlearning Protocol} (\vup) that
        accepts an Unlearning Attestation iff three \emph{independent}
        membership-inference attacks
        (shadow~\cite{shokri2017mia}, LiRA~\cite{carlini2022lira},
        label-only~\cite{choquettechoo2021labelonly}) all fall below
        $\tauMIA$, certified by $t$-of-$n$ ECDSA signatures. It is
        method-agnostic over the unlearning algorithm.
    \item A \textbf{three-contract Solidity architecture}
        (ConsentRegistry, UnlearningOracle, FLGatekeeper) deployed on
        \emph{both} a gas-free chain (PureChain, PoA$^2$) and a
        gas-based chain (Polygon Amoy, PoS), the first
        side-by-side cost comparison under both economic models.
    \item A \textbf{dual-domain empirical evaluation}: a 9-client
        N-BaIoT~\cite{meidan2018nbaiot} autoencoder federation
        (primary), cross-validated on a 5-client CASAS smart-home
        BiLSTM classifier, over three approximate-unlearning
        paradigms plus an exact-retraining baseline, measuring
        per-method MIA forgetting, on-chain gas/USD cost on both
        chains, per-round validation cost to $N\!=\!500$, and
        post-revocation utility recovery.
\end{enumerate}
